% This is samplepaper.tex, a sample chapter demonstrating the
% LLNCS macro package for Springer Computer Science proceedings;
% Version 2.21 of 2022/01/12
%
\documentclass[runningheads]{llncs}
%
\usepackage[T1]{fontenc}
% T1 fonts will be used to generate the final print and online PDFs,
% so please use T1 fonts in your manuscript whenever possible.
% Other font encondings may result in incorrect characters.
%
\usepackage{graphicx}
% Used for displaying a sample figure. If possible, figure files should
% be included in EPS format.
%
% If you use the hyperref package, please uncomment the following two lines
% to display URLs in blue roman font according to Springer's eBook style:
%\usepackage{color}
%\renewcommand\UrlFont{\color{blue}\rmfamily}
%
\usepackage{amsmath}

\begin{document}
%
\title{An algorithm for finding the global extremum 
of a partially defined function\thanks{Supported by organization x.}}
%
%\titlerunning{Abbreviated paper title}
% If the paper title is too long for the running head, you can set
% an abbreviated paper title here
%
\author{Marina Usova\inst{1}\orcidID{0000-0002-0722-6884} \and
Konstantin Barkalov\inst{1}\orcidID{0000-0001-5273-2471}}
%
\authorrunning{M. Usova et al.}
% First names are abbreviated in the running head.
% If there are more than two authors, 'et al.' is used.
%
\institute{
Department of Mathematical Software and Supercomputing Technologies, Lobachevsky University, 603022 Nizhni Novgorod, Russia
\email{\{usova,konstantin.barkalov\}@itmm.unn.ru}}
%
\maketitle              % typeset the header of the contribution
%
\begin{abstract}
The paper discusses the problem of finding the global minimum of a function that may be partially defined in the search domain. The objective function may be not defined at some search regions due to the nature of the optimized object and the simulation method (for example, numerical instability of the simulation method used). In some cases, these regions are known; but in the most cases information about them is missing. The paper gives a description of the global search algorithm for solving such class of problems. Numerical experiments confirming the efficiency of the proposed algorithm are carried out.

\keywords{Lipschitz global optimization \and multiextremal functions \and hidden constraints \and dimensionality reduction.}
\end{abstract}
%
%
%
\section{Introduction}
Please note that the first paragraph of a section or subsection is
not indented. The first paragraph that follows a table, figure,
equation etc. does not need an indent, either.

Subsequent paragraphs, however, are indented.

\section{Problem Statement}
...
\begin{equation}
\phi^*=\phi(y^* )=\min_{y \in D} \phi(y), D=\left\{ y \in R^N: a_i \leq y_i \leq b_i, a,b \in R^N,1 \leq i \leq N \right\},
\end{equation}
where $y=(y_1,y_2,...,y_N)$ - ..., $D$-dimensional hypercube, $N$ – problem dimension.
... 
$$
D=\left\{ y \in R^N: -2^{-1} \leq y_i \leq 2^{-1}, 1 \leq i \leq N \right\} = \left\{ y(x): 0 \leq x \leq 1 \right\}
$$
...
\begin{equation}
\phi^*=\phi(y(x^* ))=\min_{x \in [0,1]} \left\{ \phi(y(x)) \right\},
\end{equation}
...
\begin{equation}
| \phi (y')-\phi (y'') | \leq L \| y'-y'' \|, y',y' \in D,
\end{equation}
...
\begin{equation}
| \phi (x')-\phi (x'') | \leq K \| x'-x'' \| ^{1/N}, x',x' \in [0,1],
\end{equation}
... $K \leq 4L\sqrt N$ ...
...

\section{Method}

\subsection{Discription of ... }

\subsubsection{Global search algorithm (GSA)} ... $\phi (x)$ ... $x^1 \in (0,1)$ ... $x^{k+1}, k \leq 1$ ... $(k+1)$

\paragraph{Rule 1} Renumber (by the lower indices) the points $x^i, 1 \leq i \leq k$  ...
\begin{equation}
0=x_0 < x_1 < ... < x_i < ... < x_k=1
\end{equation}
... $z_i=\phi (x_i), 1 \leq i \leq k$ ... $v_i=v(x_i)$ ...
\begin{equation}
v_i=v(x_i)=
  \begin{cases}
    -1 & {\quad \text{if } x_i \text{ is boundary point}}\\
    1  & {\quad \text{if } x_i \text{ is internal point}}
  \end{cases}
\end{equation}
... $x_0=0$ and $x_{k+1}=1$ ... $z_0$ and $z_{k+1}$ ...

\paragraph{Rule 2} Compute current lower estimate
\begin{equation}
\mu = \max\left\{ \frac{|z_i-z_{i-1}|}{\Delta _i}, 2 \leq i \leq k \right\} , \Delta _i= \| x_i -x_{i-1} \| =(x_i-x_{i-1})^{1/N},
\end{equation}
... 
\paragraph{Rule 3} ...
\begin{equation}
z^*=\min \left\{ \phi (x_i): 1\leq i \leq k \right\}.
\end{equation}

\paragraph{Rule 4} For each interval $(x_{i-1},x_i),1 \leq i \leq k+1$ compute the value $R(i)$ called a characteristic of the interval according to the formula
\begin{equation}
R(x_i)=
  \begin{cases}
    \Delta _i+\frac {{(z_i-z_{i-1})}^2}{{(r \mu)}^2 \Delta _i} - 2 \frac {z_i+z_{i-1}-2z^*}{r \mu}, & {\quad  v(x_i)=v(x_{i-1})=1}\\
    2 \Delta _i-4 \frac {(z_i-z^*)}{r \mu}, & {\quad  v(x_{i-1})=-1, v(x_i)=1}\\
    2 \Delta _i-4 \frac {(z_{i-1}-z^*)}{r \mu}, & {\quad  v(x_{i-1})=1, v(x_i)=-1}
  \end{cases}
\end{equation}

\paragraph{Rule 5} Select interval $t$ with the largest value of characteristic $R(i):R(t)= \max_{1 \leq i \leq k+1} R(i)$.
\paragraph{Rule 6} ... $(x_{t-1},x_t)$ ...
\begin{equation}
x^{k+1}=\frac {x_t+x_{t-1}}{2}, v(x_{t-1})\neq v(x_t),
\end{equation}
...
\begin{equation}
x^{k+1}= \frac {x_t+x_{t-1}}{2} -  \text{sign} {(z_t-z_{t-1})} \frac {{|z_t-z_{t-1}|}^N}{2r \mu ^N}, v(x_{t-1})=v(x_t).
\end{equation}
...

\subsection{Imputation... }
...
\begin{itemize}
  \item Item 1
  \item Item 2
  \item Item 3
  \item Item 4
  \item Item 5
  \item Item 6
\end{itemize}
...
\begin{equation}
v_i=v(x_i)=
  \begin{cases}
    -1 & {\quad \text{if } x_i \text{ is boundary point}}\\
    1  & {\quad \text{if } x_i \text{ is internal point}}\\
    0  & {\quad \text{if } x_i \text{ is non-computability}}
  \end{cases}
\end{equation}
...

\subsubsection{First...} ...
\begin{equation}
R(i)=\Delta _i,v(x_i)=0 \text{ or } v(x_{i-1})=0
\end{equation}
... 
\begin{figure}
\includegraphics[width=\textwidth]{fig1.pdf}
\caption{Examples ... } \label{fig1}
\end{figure}
\begin{figure}
\includegraphics[width=\textwidth]{fig1.pdf}
\caption{Examples ... } \label{fig2}
\end{figure}

\subsubsection{Second...} ... 
\begin{equation}
R(i)={(1- \frac {1}{r})}^N \Delta _i,
\end{equation}
where $v(x_i)=v(x_{i-1})=0 \text{ or } v(x_i)=0 \text{ and } v(x_{i-1})=-1  \text{ or } v(x_i)=-1 \text{ and } v(x_{i-1})=0$, ...
\begin{equation}
R(i)=2\Delta _i-4 \frac {(z_i-z^*)}{r \mu}, v(x_{i-1})=-1  \text{ or } v(x_{i-1})=0, v(x_i)=1,
\end{equation}
\begin{equation}
R(i)=2\Delta _i-4 \frac {(z_{i-1}-z^*)}{r \mu}, v(x_{i-1})=1, v(x_i)=-1 \text{ or } v(x_i)=0,
\end{equation}
... on ~\ref{fig1} and ~\ref{fig2} ...

\subsubsection{Third...} ...
\begin{figure}
\includegraphics[width=\textwidth]{fig1.pdf}
\caption{Examples ... } \label{fig3}
\end{figure}
...

\subsubsection{Fourth...} ...
\begin{equation}
\tilde{z}(x)=z'+ \text{sign}(z''-z') \frac {|z''-z'| \cdot \| x-x' \|}{\| x''-x' \|},
\end{equation}
where ... 
\begin{figure}
\includegraphics[width=\textwidth]{fig1.pdf}
\caption{Examples ... } \label{fig4}
\end{figure}
\begin{figure}
\includegraphics[width=\textwidth]{fig1.pdf}
\caption{Examples ... } \label{fig5}
\end{figure}

\subsubsection{Fifth...} ...
\paragraph{Step 1} ...
\begin{equation}
\hat{x}',\hat{x}''=\frac {x'+x''}{2}\pm \text{sign}(z''-z')\cdot \frac {1}{2r} \cdot {\left(\frac {z''-z'}{\mu}\right)}^N,
\end{equation}
\paragraph{Step 2} ...
\begin{equation}
\tilde{z}(x)=
  \begin{cases}
    z', & {\quad x \leq \hat{x}',}\\
    \frac {z'+z''}{2}- \frac {r \mu}{2} (\| x - x' \| + \| x - x'' \|), & {\quad \hat{x}' < x < \hat{x}'',}\\
    z'',  & {\quad x \geq \hat{x}''.}
  \end{cases}
\end{equation}
...

\subsubsection{Discriprion of modification} ...
\begin{enumerate} 
  \item ...
  \item ...
  \item ...
\begin{equation}
z_v^*=
  \begin{cases}
    -\varepsilon _r, & {\quad v=0 ,}\\
    \min \left\{ \phi (x_i): v(x_i)=1, 1 \leq i \leq k \right\}, & {\quad v=1 ,}
  \end{cases}
\end{equation}
where $\varepsilon _r$ - ...

  \item ... $(x_{i-1}, x_i), 1 \leq i \leq k+1$ ...
      \begin{enumerate} 
        \item Case 1. ...
\begin{equation}
R(i)=
  \begin{cases}
     \Delta _i+\frac {{(z_i-z_{i-1})}^2}{{(r \mu)}^2 \Delta _i} - 2 \frac {z_i+z_{i-1}-2z_1^*}{r \mu}, & {\quad  v(x_i)=v(x_{i-1})=1}\\
    2 \Delta _i-4 \frac {(z_i-z_1^*)}{r \mu}, & {\quad  v(x_{i-1})=-1, v(x_i)=1}\\
    2 \Delta _i-4 \frac {(z_{i-1}-z_1^*)}{r \mu}, & {\quad  v(x_{i-1})=1, v(x_i)=-1}
  \end{cases}
\end{equation}
        \item Case 2. ...
\begin{equation}
R(i)=
  \begin{cases}
    2 \Delta _i-4 \frac {(z_{i+1}-z_0^*)}{r \mu}, & {\quad  v(x_{i-1})=-1, v(x_i)=0}\\
    2 \Delta _i-4 \frac {(z_{i-2}-z_0^*)}{r \mu}, & {\quad  v(x_{i-1})=0, v(x_i)=-1}
  \end{cases}
\end{equation}
        \item Case 3. ...
\begin{equation}
R(i)=
  \begin{cases}
     \Delta _i+\frac {{(z_i-z_{i-2})}^2}{{(r \mu)}^2 \Delta _i} - 2 \frac {z_i+z_{i-2}-2z_0^*}{r \mu}, & {\quad  v(x_{i-1})=0, v(x_i)=1}\\
     \Delta _i+\frac {{(z_{i+1}-z_{i-1})}^2}{{(r \mu)}^2 \Delta _i} - 2 \frac {z_{i+1}+z_{i-1}-2z_0^*}{r \mu}, & {\quad  v(x_{i-1})=1, v(x_i)=0}
  \end{cases}
\end{equation}
        \item Case 4. ...
\begin{equation}
R(i)=\Delta _i+\frac {{(\tilde{z}_i-\tilde{z}_{i-1})}^2}{{(r \mu)}^2 \Delta _i} - 2 \frac {\tilde{z}_i+\tilde{z}_{i-1}-2z_1^*}{r \mu}, v(x_{i-1})=v(x_i)=0
\end{equation}
        \item Case 5. ...
\begin{equation}
R(i)=\Delta _i \cdot {\left( 1-\frac {1}{r} \right)}^N+z_0^*.
\end{equation}
    \end{enumerate} 
  \item ...
  \item ...
\end{enumerate}

\subsection{Results... }
...  ~\ref{tab1} ...
\begin{table}
\caption{Table captions should be placed above the
tables.}\label{tab1}
\begin{tabular}{|l|l|l|l|}
\hline
 &  GKLS classical & GKLS round district & GKLS random district \\
\hline
Average number of iterations & 781,66	& 854,91 & 839,03 \\
Average number of non-computability & - & 185,79 & 42,31 \\
Number of solved problems & 100 / 100 & 100 / 100 & 100 / 100\\
\hline
\end{tabular}
\end{table}

\begin{figure}
\includegraphics[width=\textwidth]{fig1.pdf}
\caption{...} \label{fig6}
\end{figure}
...
\begin{table}
\caption{Table captions should be placed above the
tables.}\label{tab1}
\begin{tabular}{|l|l|l|l|l|}
\hline
 &  & Average number of iterations & Average number of non-computability & Solved problems \\
\hline
$N=3$ & classical & 12251,11 & - & 100 / 100 \\
$N=3$ & round district & 12201,78 & 4763,57 & 100 / 100 \\
$N=3$ & random district & 12210,71 & 63,39 & 100 / 100 \\
$N=4$ & classical & 25144,46 & - & 100 / 100 \\
$N=4$ & round district & 29840,32 & 15004,37  & 100 / 100 \\
$N=4$ & random district & 25287,49 & 26,43 & 100 / 100 \\
$N=5$ & classical & 56725,95 & - & 100 / 100 \\
$N=5$ & round district & 47840,09 & 24109,19 & 100 / 100\\
$N=5$ & random district & 57874,76 & 82,57 & 100 / 100\\
\hline
\end{tabular}
\end{table}
...
\subsection{Conclusion}
...

 (see Fig.~\ref{fig1}).

For citations of references, we prefer the use of square brackets
and consecutive numbers. Citations using labels or the author/year
convention are also acceptable. The following bibliography provides
a sample reference list with entries for journal
articles~\cite{ref_article1}, an LNCS chapter~\cite{ref_lncs1}, a
book~\cite{ref_book1}, proceedings without editors~\cite{ref_proc1},
and a homepage~\cite{ref_url1}. Multiple citations are grouped
\cite{ref_article1,ref_lncs1,ref_book1},
\cite{ref_article1,ref_book1,ref_proc1,ref_url1}.

\subsubsection{Acknowledgements} Please place your acknowledgments at
the end of the paper, preceded by an unnumbered run-in heading (i.e.
3rd-level heading).

%
% ---- Bibliography ----
%
% BibTeX users should specify bibliography style 'splncs04'.
% References will then be sorted and formatted in the correct style.
%
% \bibliographystyle{splncs04}
% \bibliography{mybibliography}
%
\begin{thebibliography}{8}
\bibitem{ref_article1}
Author, F.: Article title. Journal \textbf{2}(5), 99--110 (2016)

\bibitem{ref_lncs1}
Author, F., Author, S.: Title of a proceedings paper. In: Editor,
F., Editor, S. (eds.) CONFERENCE 2016, LNCS, vol. 9999, pp. 1--13.
Springer, Heidelberg (2016). \doi{10.10007/1234567890}

\bibitem{ref_book1}
Author, F., Author, S., Author, T.: Book title. 2nd edn. Publisher,
Location (1999)

\bibitem{ref_proc1}
Author, A.-B.: Contribution title. In: 9th International Proceedings
on Proceedings, pp. 1--2. Publisher, Location (2010)

\bibitem{ref_url1}
LNCS Homepage, \url{http://www.springer.com/lncs}. Last accessed 4
Oct 2017
\end{thebibliography}
\end{document}
